\scanpage{131}
\mspage{16}
$\tau M=(\tau x_1,\ldots,\tau x_n)(\tau R)$ and $\tau x_1\ne0$ and hence there exists a unique positive integer $s$ such that $(\tau x_1)^{\unsure{S|s}}\in\tau N$ and $(\tau x_1)^{\unsure{S|s}+1}\notin\tau N$ and therefore as in the proof of [C\deleted{\unsure{0|O}}\inserted{o}, Theorem 8],
\margin{small ``oh''}
$y=y'\eta$ where $y'=x_1^{\unsure{S|s}}+y_1x_1^{\unsure{S|s}-1}+\cdots+y_s$ with $y_i\in N$ and $\eta$ a unit in $R$. Since the \inserted{line} $x_2=\cdots=x_n=0$ is non-tangential to $y$ it is also non-tangential to $y'$ and hence $\Lambda(y')$ does not vanish at $x_2=\cdots=x\deleted{n}\inserted{{}_n}=0$ and from this, in view of the fact that $y_i\in N$, we can conclude that $s=\lambda_R(y')=\lambda_R(y)=b$.

To prove the uniqueness of $y'$ abandon the assumption that $y$ is an irreducible not divisible by $x_1$ and let $y_1^*,\ldots,y_b^*$ be elements in $S$ and $\eta^*$ a unit in $R$ such that $y=y^*\eta^*$ with $y^*=x_1^b+y_1^*x_1^{b-1}+\cdots+y_1^*$. Then $y'=y^*(\eta^*/\eta)$ where $(\eta^*/\eta)$ is a unit in $R$, $\lambda_R(y^*)=\lambda_R(y')=b$, and the coefficient of $x_1^b$ in $y^*$ as well as $y'$ is $1$, and hence $\Lambda(y')=\Lambda(y^*)$ and $\Lambda(\eta^*/\eta)=1$. Now $y'=\sum_{i=0}^\infty\alpha_{b+\deleted{1}\inserted{i}}$, $y^*=\sum_{i=0}^\infty\beta_{b+i}$, $\eta^*/\eta=1+\sum_{i=1}^\infty\gamma_{\deleted{1}\inserted{i}}$ where $\alpha_b$ and $\beta_b$ are (non-zero) forms of degree $b$ in $k[x_1,\ldots,x_n]$ and have $x_1$-degree $b$; for $i>0$, $\alpha_{b+i}$ is either zero or is a form of degree $b+i$ in $k[x_1,\ldots,x_n]$ and has $x_1$-degree less than $b$, $\beta_{b+i}$ is either zero or is a form of degree $b+i$ in $k[x_1,\ldots,x_n]$ and has $x_1$-degree less than $b$ and $\gamma_i$ is either zero or is a form of degree $i$ in $k[x_1,\ldots,x_n]$. By induction on $i$ we shall show that $\gamma_i=0$ for all $i>0$ and this will show that $y'=y^*$ and the uniqueness statement will have been proved. Now $\alpha_{b+1}=\beta_b\gamma_1+\beta_{b+1}$, i.e., $\beta_b\gamma_1=\alpha_{b+1}-\beta_{b+1}$, since the $x_1$-degree of $\beta_b$ is $b$, $\alpha_{b+1}-\beta_{b+1}$ is either zero has $x_1$-degree less than $b$, hence
\mspage{130}
